% संपूर्ण मराठी ग्रंथातील प्रकरण 38: लॅम्डा-परिभाष्यता
% हा संपादनयोग्य स्रोतखंड आहे. संकलनासाठी संपूर्ण स्रोतसंग्रहातील
% अक्षररूपे, पूर्वभाग, चिन्हव्याख्या आणि आकृती-स्रोत आवश्यक आहेत.
% मूळ: Open Logic Project; मजकूर आणि रूपांतर CC BY 4.0.
% यंत्रानुवाद, दुरुस्ती आणि तपासणी: OpenAI Codex —
% GPT-5.6 Sol आणि GPT-6 Sol, दोन्ही Ultra effort.
% दुरुस्ती-पुनरावलोकन: GPT-6.1 Sol, Ultra effort.
\chapter{लॅम्डा-परिभाष्यता}\label{lam:rep:chap}
\begin{quote}
हे प्रकरण प्रायोगिक स्वरूपाचे आहे. अधिक स्पष्टीकरणाची
  गरज आहे; व्याख्या, सिद्धतांसह विधाने आणि अधिक
  उदाहरणे अशा रचनेत मजकूर अधिक चांगल्या प्रकारे
  मांडायला हवा.
\end{quote}








\section{परिचय}\label{lam:ldf:int:sec}

पहिल्या नजरेला लॅम्डा कलन हे फलने दर्शवणाऱ्या
अभिव्यक्ती आणि त्यांचे इतर अभिव्यक्तींना उपयोजन
यांचे अतिशय अमूर्त कलन वाटते. या कलनाच्या
विन्यासात ही फलने कोणत्या विशिष्ट प्रकारच्या वस्तूंवर
असतात याचा निर्देश नाही; विशेषतः नैसर्गिक संख्यांचा
उल्लेखही नाही. त्याच्या मूलभूत क्रिया—उपयोजन आणि
लॅम्डा अमूर्तीकरण—सर्व प्रकारच्या फलनांवर लागू
होतात, केवळ नैसर्गिक संख्यांवरील फलनांवर नाही.

तरीही थोड्या कल्पकतेने लॅम्डा कलनात अंकगणितीय
फलने, म्हणजे नैसर्गिक संख्यांवरील फलने, परिभाषित
करता येतात. त्यासाठी प्रत्येक नैसर्गिक संख्या~$n \in
\Nat$ करिता विशेष $\lambd$-पद~$\num n$ परिभाषित
करू. ते~$n$ साठीचे \emph{चर्च अंक} आहे. चर्च अंकांना
अलाँझो चर्च यांचे नाव दिले आहे.

\begin{defn}
  $n \in \Nat$ असेल, तर संबंधित \emph{चर्च अंक}
  $\num{n}$ हे~$n$ दर्शवते:
  \[    \num{n} \ident \lambd[fx][f^n(x)]
  \]
  येथे $f^n(x)$ म्हणजे $f$ हे~$x$ वर $n$ वेळा लागू
  केल्याचा परिणाम. उदाहरणार्थ, $\num{0}$ हे
  $\lambd[fx][x]$ आहे आणि $\num{3}$ हे
  $\lambd[fx][f(f(f\,x))]$ आहे.
\end{defn}

चर्च अंक $\num n$ हे दोन फलसाधक $f$ आणि $x$
स्वीकारून $f^n(x)$ परत करणारे फलन दर्शवणारे लॅम्डा
पद म्हणून संकेतित केलेले आहे. चर्च अंक स्पष्टपणे
सामान्य रूपात आहेत.

लॅम्डा कलनात नैसर्गिक संख्या दर्शवणे उपयोगी ठरते
तेव्हाच, जेव्हा त्यांच्यावर संगणनही करता येते.
चर्च अंकांवर लॅम्डा कलनात संगणन करणे म्हणजे
एखादे $\lambd$-पद~$F$ अशा चर्च अंकाला लागू करणे
आणि तयार झालेले पद~$F\, \num n$ सामान्य रूपात
न्यूनीत करणे. ते नेहमी सामान्य रूपात न्यूनीत झाले
आणि ते सामान्य रूप नेहमी एखादे चर्च अंक~$\num m$
असेल, तर संगणनाचा परिणाम~$m$ ही संख्या मानता येते.
मग~$F$ हे $f\colon \Nat \to \Nat$ फलन परिभाषित
करते असे मानता येते: $f(n) = m$ असते, जर आणि फक्त जर
$F\, \num n \red \num m$. चर्च–रॉसर गुणधर्मामुळे
सामान्य रूप अस्तित्वात असेल, तर ते एकमेव असते.
म्हणून $F\, \num n \red \num m$ असेल, तर
$F \, \num n$ चे दुसऱ्या कोणत्याही सामान्य रूपात,
विशेषतः दुसऱ्या चर्च अंकात, न्यूनीकरण होऊ शकत नाही.

उलट, $f\colon \Nat \to \Nat$ फलन दिले असता,
अशा रीतीने~$f$ परिभाषित करणारे पद $F$ आहे का,
असे विचारता येते. असे असेल, तर $F$ हे~$f$ ला
\emph{लॅम्डा-परिभाषित} असे म्हणतो आणि $f$ हे
लॅम्डा-परिभाष्य आहे असे म्हणतो. ही संकल्पना
अनेकस्थानी आणि आंशिक फलनांपर्यंत व्यापक करता येते.

\begin{defn}
$f\colon \Nat^k \to \Nat$ असे मानू; येथे आंशिक
फलनही मान्य आहे. या प्रकरणातील न्यूनीकरण आणि सामान्य
रूप हे बीटा-न्यूनीकरण आणि बीटा-सामान्य रूप या अर्थाने आहेत. सर्व
$n_0$, \dots,~$n_{k-1}$ साठी $f(n_0, \dots,
n_{k-1})$ परिभाषित असेल, तर
\[F \, \num{n_0} \, \num{n_1} \dots \num{n_{k-1}} \red \num{f(n_0, n_1, \dots,
  n_{k-1})}
\]
आणि अन्यथा $F \, \num{n_0} \,
\num{n_1} \dots \num{n_{k-1}}$ ला सामान्य रूप
नसेल, तेव्हा बंद लॅम्डा पद~$F$ हे $f$ ला
\emph{लॅम्डा-परिभाषित} असे म्हणतो.
\end{defn}

स्थिर फलने हे एक अत्यंत सोपे उदाहरण आहे.
$C_k \ident \lambd[x][\num{k}]$ हे पद
$c_k\colon \Nat \to \Nat$ फलनाला
लॅम्डा-परिभाषित करते; येथे $c_k(n) = k$.
कारण कोणत्याही~$n$ साठी $C_k \, \num n \ident
(\lambd[x][\num{k}])\num n \redone \num{k}$.
ओळख फलन $\lambd[x][x]$ ने लॅम्डा-परिभाषित आहे.
अधिक गुंतागुंतीची फलने परिभाषित करणे अर्थातच
कठीण असते आणि त्यासाठी बरीच कल्पकता लागते.
त्यामुळे प्रत्येक संगणनक्षम फलन लॅम्डा-परिभाष्य
आहे हे काहीसे आश्चर्यकारक वाटू शकते. याचा उलट
दिशेचा निष्कर्षही खरा आहे: फलन लॅम्डा-परिभाष्य
असेल, तर ते संगणनक्षम असते.











\section{लॅम्डा-परिभाष्य अंकगणितीय फलने}\label{lam:rep:arf:sec}

\begin{prop}
  \label{lam:rep:arf:prop:succ-ld}
  उत्तरवर्ती फलन~$\Succ$ हे लॅम्डा-परिभाष्य आहे.
\end{prop}

\begin{proof}
उत्तरवर्ती फलनाला लॅम्डा-परिभाषित करणारे पद असे आहे:
\begin{align*}
  \fn{Succ} & \ident \lambd[a][\lambd[fx][f ({a} f x)]].
  \intertext{आपल्या संकेतनरूढीनुसार याचा पूर्ण अर्थ असा:}
  \fn{Succ} & \ident \lambd[a][\lambd[f][\lambd[x][(f (({a} f) x))]]].
  \intertext{$\fn{Succ}$ हे~$a$ हा संख्यारूपी फलसाधक
    स्वीकारणारे फलन आहे आणि त्याचे मूल्य दुसरे फलन,
    $\lambd[fx][f ({a} f x)]$, आहे. ते फलन स्वतः चर्च अंक
    नाही. मात्र $a$ हा फलसाधक चर्च अंक असेल, तर त्याचे
    न्यूनीकरण चर्च अंकात होते. पाहा:}
   (\lambd[a][\lambd[fx][f ({a} f x)]])\,\num{n} & \redone
   \lambd[fx][f ({\num{n}} f x)].
   \intertext{आतील पद $\num{n}fx$ मध्ये रेडेक्स आहे, कारण
     $\num{n}$ हे $\lambd[fx][f^nx]$ आहे. म्हणून
     $\num{n}fx \red f^nx$. त्यामुळे संपूर्ण पदासाठी}
   \fn{Succ}\, \num n & \red \lambd[fx][f(f^n(x))],
\end{align*}
म्हणजे $\num{n+1}$.
\end{proof}

\begin{ex}
  $\fn{Succ}$ हे~$\num{0}$, म्हणजे
  $\lambd[fx][x]$, ला लागू केल्यावर काय होते ते पाहू.
  पदे पूर्णपणे लिहू:
  \begin{align*}
    \fn{Succ}\, \num{0} & \ident (\lambd[a][\lambd[f][\lambd[x][(f (({a} f) x))]]])(\lambd[f][\lambd[x][x]])\\
    & \redone \lambd[f][\lambd[x][(f (({(\lambd[f][\lambd[x][x]])} f) x))]] \\
    & \redone \lambd[f][\lambd[x][(f ({(\lambd[x][x])} x))]] \\
    & \redone \lambd[f][\lambd[x][(f x)]] \ident \num{1}\\
  \end{align*}
\end{ex}

\begin{prob}
  पुढील पद
  \begin{align*}
    \fn{Succ}' & \ident \lambd[{n}][\lambd[fx][{n} f (f x)]]
  \end{align*}
  उत्तरवर्ती फलनाला लॅम्डा-परिभाषित करते. का ते स्पष्ट करा.
\end{prob}

\begin{prop}
  \label{lam:rep:arf:prop:add-ld}
  बेरीज फलन~$\Add$ हे लॅम्डा-परिभाष्य आहे.
\end{prop}

\begin{proof}
बेरीज पुढील पदांनी लॅम्डा-परिभाषित होते:
\begin{align*}
  \fn{Add} & \ident \lambd[{a}{b}][\lambd[fx][{a} f ({b} f x)]]
\intertext{किंवा पर्यायाने,}
  \fn{Add}' & \ident \lambd[{a}{b}][{a}\, \fn{Succ} \, {b}].
\intertext{पहिले बेरीजपद असे कार्य करते: $\fn{Add}$ प्रथम
  ${a}$ आणि ${b}$ या दोन संख्या स्वीकारते. त्यातून
  $f$ आणि $x$ स्वीकारून $af(bfx)$ परत करणारे फलन
  मिळते. $a$ आणि $b$ हे चर्च अंक $\num{n}$ आणि
  $\num{m}$ असतील, तर त्याचे न्यूनीकरण $f^{n+m}(x)$
  मध्ये होते; ते $f^{n}(f^{m}(x))$ सारखेच आहे.
  टप्प्याटप्प्याने पाहू:}
  (\lambd[{a}{b}][\lambd[fx][{a} f ({b} f x)]])\num n\,\num m & \red
  \lambd[fx][\num{n}\, f (\num {m}\, f x)] \\
  & \red \lambd[fx][\num{n}\, f (f^m x)] \\
  & \red \lambd[fx][f^n (f^m x)] \ident \num {n+m}.
\intertext{बेरीजचे दुसरे निरूपण $\fn{Add'}$ वेगळे कार्य
  करते: ते चर्च अंक $\num{n}$ आणि $\num{m}$ यांना
  लागू केल्यास,}
\fn{Add}' \num n \,\num m
& \red \num{n}\, \fn{Succ}\, \num{m}.
\intertext{पण $\num{n} f x$ चे नेहमी $f^n(x)$ मध्ये
  न्यूनीकरण होते. म्हणून,}
  \num{n}\,  \fn{Succ}\, \num{m} & \red \fn{Succ}^n(\num{m}).
\end{align*}
$\fn{Succ}$ हे उत्तरवर्ती फलनाला लॅम्डा-परिभाषित
करते आणि~$m$ ला उत्तरवर्ती फलन $n$ वेळा लागू
केल्यावर $n+m$ मिळते. म्हणून याचे पुढे
$\num{n+m}$ मध्ये न्यूनीकरण होते.
\end{proof}

\begin{prop}
  \label{lam:rep:arf:prop:mult-ld}
  गुणाकार पुढील पदाने लॅम्डा-परिभाष्य आहे:
  \[  \fn{Mult} \ident \lambd[ab][\lambd[fx][a (b f) x]]
  \]
\end{prop}

\begin{proof}
  हे कसे कार्य करते ते पाहण्यासाठी $\fn{Mult}$ ला
  चर्च अंक $\num{n}$ आणि $\num{m}$ लागू करू.
  $\fn{Mult} \, \num{n} \, \num{m}$ चे न्यूनीकरण
  $\lambd[fx][\num{n}(\num{m}\, f)x]$ मध्ये होते.
  पद $\num{m} f$ हे आपल्या फलसाधकाला $f$ हे
  $m$ वेळा लागू करणारे फलन परिभाषित करते.
  परिणामी, $\num{n} (\num{m} f) x$ हे
  ``$f$ हे $m$ वेळा लागू करा'' या फलनाला~$x$ वर
  $n$ वेळा लागू करते. म्हणजे $f$ हे $x$ वर
  $n\cdot m$ वेळा लागू होते. मिळणारे सामान्य पद
  म्हणजे नेमका चर्च अंक $\num{nm}$.
\end{proof}

\begin{editorial}
$\eta$-न्यूनीकरणाने हे पद आणखी सोपे करता येते:
\[  \fn{Mult} \ident \lambd[ab][\lambd[f][a (b f)]].
\]

पण त्याआधी $\eta$-न्यूनीकरण स्पष्ट करावे लागेल.
\end{editorial}

\begin{prob}
पुढील पदाने गुणाकार लॅम्डा-परिभाषित होतो:
\[  \fn{Mult}' \ident \lambd[ab][b (\fn{Add}\, a) \num{0}].
\]
हे का कार्य करते ते स्पष्ट करा.
\end{prob}


घातांक फलनाची $\lambd$-पद म्हणून व्याख्या
आश्चर्यकारकपणे सोपी आहे:
\[  \fn{Exp} \ident \lambd[be][\lambd[fx][e b f x]].
\]
पहिला फलसाधक $b$ हा पाया आणि दुसरा~$e$ हा
घातांक आहे. आपल्या संख्या-संकेतनानुसार $e f$
हे सहजपणे $f^e$ आहे. मात्र चर्च अंकाचे पूर्ण
बीटा-सामान्य रूप मिळण्यासाठी दोन बाह्य प्राचल आवश्यक
आहेत. घातांक शून्य असल्यास आतले पद identity मध्ये
न्यूनीत होते आणि या दोन प्राचलांनी चर्च अंक एक मिळतो.
इतर घातांकांसाठी पद वारंवार फलनसंयोजन करून योग्य
चर्च अंक देते. येथे शून्याच्या शून्य घाताचे मूल्यही एक
मानले आहे. हे समजणे कठीण वाटल्यास,
वारंवार गुणाकारानेही घातांक फलन परिभाषित करता येते:
\[  \fn{Exp}' \ident \lambd[be][e (\fn{Mult}\, b) \num{1}].
\]

चर्च अंकांवरील पूर्ववर्ती फलन आणि वजाबाकी वाटतात
तितकी सोपी नाहीत: त्यांसाठी जोड्यांचे संकेतन आवश्यक आहे.











\section{जोड्या आणि पूर्ववर्ती फलन}\label{lam:ldf:pai:sec}

\begin{defn}
$M$ आणि $N$ यांची जोडी ($\tuple{M,N}$ असे लिहितात)
पुढीलप्रमाणे परिभाषित केली जाते. येथे अमूर्तीकरणाचे
नाव दोन्ही पदांच्या मुक्त चरांपासून वेगळे निवडायचे आहे:
\[\tuple{M,N} \ident \lambd[f][fMN].
\]
\end{defn}

सहज अर्थाने ही जोडी म्हणजे एक फलन: ते दुसरे फलन
स्वीकारते आणि ते जोडीतल्या दोन्ही घटकांना लागू करते.
या कल्पनेवरून दोन पदे घेऊन त्यांची जोडी परत करणारे
पुढील रचनाफलन मिळते:
\[  \fn{Pair} \ident \lambd[mn][\lambd[f][fmn]]
\]
जोडी दिल्यावर तिचे घटक पुन्हा मिळवायचे असतात.
त्यासाठी जोडी फलसाधक म्हणून स्वीकारून तिचा पहिला
किंवा दुसरा घटक परत करणारी दोन घटक-निवड फलने हवीत:
\begin{align*}
  \fn{Fst} & \ident \lambd[p][p(\lambd[mn][m])]\\
  \fn{Snd} & \ident \lambd[p][p(\lambd[mn][n])]
\end{align*}

\begin{prob}
  घटक-निवड फलने $\fn{Fst}$ आणि $\fn{Snd}$ का कार्य
  करतात ते स्पष्ट करा.
\end{prob}

आता जोड्यांच्या साहाय्याने पूर्ववर्ती फलन
लॅम्डा-परिभाषित करता येते:
\[  \fn{Pred} \ident \lambd[n][\fn{Fst}(n (\lambd[p][\tuple{\fn{Snd}\, {p}, \fn{Succ}(\fn{Snd}\, {p})}]) \tuple{\num 0, \num 0})]
\]
$\num n\, f x$ चे $f^{n}(x)$ मध्ये न्यूनीकरण
होते हे लक्षात ठेवा. येथे $f$ हे जोडी $p$ स्वीकारून
$p$ चा दुसरा घटक आणि त्या दुसऱ्या घटकाचा
उत्तरवर्ती यांची नवी जोडी परत करणारे फलन आहे;
$x$ ही $\tuple{\num{0},\num{0}}$ जोडी आहे. म्हणून $n=0$
असेल, तर परिणामी जोडी $\tuple{\num{0},\num{0}}$ असते;
अन्यथा ती $\tuple{\num{n-1}, \num n}$ असते.
मग $\fn{Pred}$ या परिणामी जोडीचा पहिला घटक
परत करते.

नैसर्गिक संख्यांवरील शून्यावर थांबणारी वजाबाकी
म्हणजे $a$ ला $\fn{Pred}$ हे $b$ वेळा लागू करणे:
\[\fn{Sub} \ident \lambd[ab][b \fn{Pred}\, a].
\]











\section{सत्यतामूल्ये आणि संबंध}\label{lam:ldf:tvr:sec}

शुद्ध लॅम्डा कलनात सत्यतामूल्यांचे पुढीलप्रमाणे
संकेतन करता येते:
\begin{align*}
  \fn{true} & \ident \lambd[x][\lambd[y][x]]\\
  \fn{false} & \ident \lambd[x][\lambd[y][y]]
\end{align*}

सत्यतामूल्ये \emph{निवडफलने} म्हणून दर्शवली जातात:
दोन फलसाधक स्वीकारून त्यांपैकी एक परत करणारी
फलने. सत्यतामूल्य $\fn{true}$ पहिला फलसाधक
निवडते आणि $\fn{false}$ दुसरा. उदाहरणार्थ,
$\fn{true}\, M N$ चे नेहमी $M$ मध्ये, तर
$\fn{false}\, M N$ चे नेहमी~$N$ मध्ये न्यूनीकरण होते.

\begin{defn}
संबंध $R \subseteq \Nat^k$ हा लॅम्डा-परिभाष्य
आहे असे म्हणतो, जर आणि फक्त जर असे बंद पद~$R$ असते की
प्रत्येक नैसर्गिक निविष्टांसाठी पुढील अटी पूर्ण होतात:
\begin{align*}
  R\, \num{n_1} \dots \num{n_k} & \bred \fn{true}
  \intertext{जेव्हा $R(n_1, \dots, n_k)$ सत्य असेल, आणि}
  R\, \num{n_1} \dots \num{n_k} & \bred \fn{false}
\end{align*}
अन्यथा.
\end{defn}

उदाहरणार्थ, $0$ साठी खरा आणि $0$ व्यतिरिक्त इतर
कोणत्याही संख्येसाठी खरा नसलेला संबंध
$\fn{IsZero} = \{0\}$ पुढील पदाने
लॅम्डा-परिभाष्य आहे:
\[  \fn{IsZero} \ident \lambd[n][n (\lambd[x][\fn{false}])\, \fn{true}].
\]
हे कसे कार्य करते? चर्च अंक हे पुनरावर्तक आहेत:
पहिला फलसाधक दुसऱ्यावर $n$~वेळा लागू करणारी
फलने. म्हणून सुरुवातीचे मूल्य $\fn{true}$ घेऊ;
पुनरावर्तनाच्या प्रत्येक पायरीत आधीच्या पायरीचा
निकाल काहीही असला, तरी $\fn{false}$ परत करू.
ही पायरी सुरुवातीच्या मूल्यावर $n$ वेळा लागू होते.
म्हणून ती एकदाही लागू झाली नाही, म्हणजे $n = 0$
असेल, तेव्हाच निकाल $\fn{true}$ असतो.

सत्यतामूल्यांच्या या संकेतनावरून आणखी सत्यता-फलने
परिभाषित करता येतात. नकरण आणि संधियोग यांची
दोन निरूपणे अशी:
\begin{align*}
  \fn{Not} & \ident \lambd[x][x\, \fn{false}\, \fn{true}]\\
  \fn{And} & \ident \lambd[x][\lambd[y][xy \,\fn{false}]]
\end{align*}
``$\fn{Not}$'' हे एक फलसाधक स्वीकारते:
तो $\fn{false}$ असेल, तर $\fn{true}$ परत करते;
आणि तो~$\fn{true}$ असेल, तर $\fn{false}$.
``$\fn{And}$'' हे दोन सत्यतामूल्ये स्वीकारते आणि
निकाल $\fn{true}$ असायला हवा, जर आणि फक्त जर
दोन्ही निविष्ट सत्यतामूल्ये~$\fn{true}$ असतील. वर सांगितल्याप्रमाणे सत्यतामूल्ये
निवडफलने आहेत. म्हणून सत्यतामूल्य $x$ ला दोन
फलसाधक लागू केल्यास, $x$ हे $\fn{true}$ असेल
तर पहिला फलसाधक आणि अन्यथा दुसरा मिळतो.
आता $\fn{And}$ आपले दोन फलसाधक $x$ आणि $y$
घेते व $x$ ला अनुक्रमे $y$ आणि $\fn{false}$
हे देते. $x$ सत्यतामूल्य आहे असे गृहीत धरल्यास,
$x$ हे $\fn{true}$ असेल तर निकाल~$y$,
आणि $x$ हे $\fn{false}$ असेल तर निकाल
$\fn{false}$ होतो; हाच अपेक्षित परिणाम आहे.

येथे $\fn{And}$ ला फलसाधक म्हणून फक्त
सत्यतामूल्येच दिली जातात असे आपण गृहीत धरतो.
इतर पदे दिल्यास निकाल—सामान्य रूप अस्तित्वात
असेल तर ते—सत्यतामूल्य असेलच असे नाही.

\begin{prob}
  सत्यतामूल्यांचे $\lambd$-पदांद्वारे केलेले संकेतन
  वापरून समावेशक आणि अपवर्जक विकल्पयोग दर्शवणारी
  $\fn{Or}$ आणि $\fn{Xor}$ ही फलने परिभाषित करा.
\end{prob}











\section{आदिम पुनरावर्ती फलने लॅम्डा-परिभाष्य असतात}\label{lam:ldf:prf:sec}

मूलभूत फलने $\Zero$, $\Succ$ आणि $\Proj{n}{i}$
यांपासून संयोजन व आदिम पुनरावर्तन वापरून परिभाषित
करता येणारी फलने म्हणजे आदिम पुनरावर्ती फलने, हे लक्षात घ्या.

\begin{lem}
  \label{lam:ldf:prf:lem:basic}
  मूलभूत आदिम पुनरावर्ती फलने $\Zero$, $\Succ$ आणि
  प्रक्षेपणे~$\Proj{n}{i}$ ही लॅम्डा-परिभाष्य आहेत.
\end{lem}

\begin{proof}
पुढील पदे त्यांना लॅम्डा-परिभाषित करतात:
\begin{align*}
  \fn{Zero} & \ident \lambd[a][\lambd[fx][x]]\\
  \fn{Succ} & \ident \lambd[a][\lambd[fx][f (a f x)]] \\
  \fn{Proj}^n_i & \ident \lambd[x_0\dots x_{n-1}][x_i]
\end{align*}
\end{proof}

\begin{lem}
  \label{lam:ldf:prf:lem:comp} सर्वत्र परिभाषित $k$-स्थानी फलन $f$ आणि $n$-स्थानी
  सर्वत्र परिभाषित फलने $g_0, \dots, g_{k-1}$ ही अनुक्रमे $F$, $G_0$,
  \dots, $G_{k-1}$ या पदांनी लॅम्डा-परिभाष्य आहेत,
  आणि $h$ हे त्यांच्यापासून संयोजनाने परिभाषित आहे असे समजा.
  मग $h$ हे लॅम्डा-परिभाष्य आहे.
\end{lem}

\begin{proof}
  पुढील पद $h$ ला लॅम्डा-परिभाषित करते:
  \[  H \ident \lambd[x_0 \dots x_{n-1}][F\, (G_0 x_0 \dots
    x_{n-1}) \dots (G_{k-1} x_0 \dots x_{n-1})]
  \]
  याची पडताळणी सरावासाठी सोडली आहे.
\end{proof}

\begin{prob}
  $H\num{n_0}\dots\num{n_{n-1}} \red \num{h(n_0, \dots,
    n_{n-1})}$ हे दाखवून \ref{lam:ldf:prf:lem:comp}
  चा पुरावा पूर्ण करा.
\end{prob}

\ref{lam:ldf:prf:lem:comp} मध्ये $f$ आणि $g_0$,
\dots,~$g_{k-1}$ ही आदिम पुनरावर्ती असणे आवश्यक
नव्हते, हे लक्षात घ्या; ती सर्वत्र परिभाषित आणि
लॅम्डा-परिभाष्य असणे पुरेसे आहे.

\begin{lem}\label{lam:ldf:prf:lem:prim}
  $f$ हे $n$-स्थानी आणि~$g$ हे $n+2$-स्थानी फलन
  असून ती सर्वत्र परिभाषित आणि अनुक्रमे $F$ आणि~$G$ या पदांनी
  लॅम्डा-परिभाष्य आहेत असे समजा. $h$ हे
  $f$ आणि $g$ यांच्यापासून आदिम पुनरावर्तनाने परिभाषित असेल,
  तर $h$ देखील लॅम्डा-परिभाष्य आहे.
\end{lem}

\begin{proof}
  $h$ ची व्याख्या अशी आहे, हे लक्षात घ्या:
  \begin{align*}
    h(x_1, \dots, x_n, 0) &= f(x_1, \dots, x_n)\\
    h(x_1, \dots, x_n, y+1) & = g(x_1, \dots, x_n, y, h(x_1, \dots, x_n, y)).
  \end{align*}
  सहज अर्थाने, आदिम पुनरावर्ती व्याख्येत $f(x_1,
  \dots, x_n)$ या सुरुवातीच्या मूल्यापासून $h$ चे
  मूल्य $y$ वेळा अद्ययावत होते. हे चर्च अंकांच्या
  व्याख्येसारखे आहे: तेही पुनरावर्तक आहेत.

  सोपेपणासाठी आपण एका अतिरिक्त फलसाधक~$x$ साठी
  व्याख्या व पुरावा देऊ. $h$ ला लॅम्डा-परिभाषित
  करणारे पद असे आहे:
  \begin{align*}
    H \ident & \lambd[x][\lambd[y][\fn{Snd} (y
        D \tuple{\num{0}, F x})]]
    \intertext{जेथे }
    D \ident & \lambd[p][\tuple{\fn{Succ} (\fn{Fst}\, p), (G x
          (\fn{Fst}\, p) (\fn{Snd}\, p))}]
  \end{align*}
  पुनरावर्तनात आपण ठेवलेली अवस्था ही जोडी आहे:
  पहिला घटक सध्याचे $y$ आणि दुसरा त्या स्थितीसाठीचे
  $h$ चे मूल्य. प्रत्येक पायरीत $y$ आणि $h$
  यांच्या नव्या मूल्यांची जोडी तयार होते. पुनरावर्तन
  संपल्यावर तिचा दुसरा घटक परत करतो; तेच $h$ चे
  अंतिम मूल्य असते. एकाहून अधिक अतिरिक्त फलसाधक
  असतील, तर ती सर्व निविष्टे सुरुवातीच्या फलनाला आणि
  प्रत्येक पायरीच्या फलनाला त्याच क्रमाने दिली जातात;
  प्रत्येक बाह्य प्राचलासाठी स्वतंत्र अमूर्तीकरण घ्यायचे.
  पुनरावर्तनाच्या जोडीतील मोजणी आणि पूर्वमूल्याचे घटक
  तेच राहतात, म्हणून खालील विगमन तसाच लागू होतो.
  हे निरूपण आदिम पुनरावर्तनाचेच
  आहे, हे आता सिद्ध करू.

  कोणत्याही $n$ आणि $m$ साठी $H\,\num{n}\,\num{m} \red
  \num{h(n,m)}$ दाखवायचे आहे. त्यासाठी आधी $D_n \ident
  \Subst{D}{\num{n}}{x}$ असेल, तर $D_n^m \tuple{\num{0}, F\, \num{n}} \red
  \tuple{\num{m}, \num{h(n, m)}}$ हे दाखवू. $m$ वरील
  विगमनाने पुढे जाऊ.

  $m=0$ असेल, तर $D_n^0 \tuple{\num{0}, F\, \num{n}} \red
  \tuple{\num{0}, \num{h(n, 0)}}$ दाखवायचे आहे.
  पण $D_n^0 \tuple{\num{0}, F\,
    \num{n}}$ म्हणजेच $\tuple{\num{0}, F\, \num{n}}$.
  $F$ हे $f$ ला लॅम्डा-परिभाषित करते; त्यामुळे याचे
  $\tuple{\num{0}, \num{f(n)}}$ मध्ये न्यूनीकरण होते;
  आणि $f(n) = h(n, 0)$ असल्याने हेच
  $\tuple{\num{0}, \num{h(n,0)}}$ आहे.

  आता $D_n^m \tuple{\num{0}, F\, \num{n}} \red
  \tuple{\num{m}, \num{h(n, m)}}$ असे समजा.
  $D_n^{m+1} \tuple{\num{0}, F\, \num{n}} \red
  \tuple{\num{m+1}, \num{h(n, m+1)}}$ हे दाखवू:
  \begin{align*}
    D_n^{m+1} \tuple{\num{0}, F\, \num{n}}
    & \ident D_n(D_n^{m} \tuple{\num{0}, F\, \num{n}})\\
    & \red D_n\,\tuple{\num{m}, \num{h(n, m)}} \text{\qquad (विगमन गृहीतकाने)}\\
    & \ident (\lambd[p][
      \tuple{
        \fn{Succ} (\fn{Fst}\, p), (G \, \num{n} (\fn{Fst}\, p) (\fn{Snd}\, p))
    }]) \tuple{\num{m}, \num{h(n,m)}}\\
    & \redone
    \tuple{\fn{Succ} (\fn{Fst}\, \tuple{\num{m}, \num{h(n,m)}}),\\
      & \qquad (G \, \num{n} (\fn{Fst}\, \tuple{\num{m}, \num{h(n,m)}}) (\fn{Snd}\, \tuple{\num{m}, \num{h(n,m)}}))}\\
    & \red
    \tuple{\fn{Succ}\, \num{m}, (G \, \num{n} \,\num{m} \, \num{h(n,m)})}\\
    & \red \tuple{\num{m+1}, \num{g(n, m, h(n, m))}}
  \end{align*}
  $g(n, m, h(n, m)) = h(n, m+1)$ असल्याने हे सिद्ध झाले.

  शेवटी पुढील न्यूनीकरण पाहा:
  \begin{align*}
    H\,\num n\,\num m &\ident \lambd[x][\lambd[y][\fn{Snd} (y
        (\lambd[p].\tuple{\fn{Succ} (\fn{Fst}\, p), (G\, x\,
          (\fn{Fst}\, p)\, (\fn{Snd}\, p))}) \tuple{\num{0}, F x})]]\\
    & \qquad\,\num n\, \num m\\
    & \red \fn{Snd} (\num m \,
    \underbrace{(\lambd[p].\tuple{\fn{Succ} (\fn{Fst}\, p), (G \,\num n\,
      (\fn{Fst}\, p) (\fn{Snd}\, p))})}_{D_n} \tuple{\num{0}, F \num n})\\
    & \ident \fn{Snd} (\num m \, D_n\, \tuple{\num{0}, F \num n})\\
    & \red \fn{Snd} \,(D_n^m  \tuple{\num{0}, F \num n})\\
    & \red \fn{Snd} \,\tuple{\num{m}, \num{h(n, m)}} \\
    & \red \num{h(n,m)}.
  \end{align*}
\end{proof}

\begin{prop}
  प्रत्येक आदिम पुनरावर्ती फलन लॅम्डा-परिभाष्य आहे.
\end{prop}

\begin{proof}
  \ref{lam:ldf:prf:lem:basic} नुसार सर्व मूलभूत फलने
  लॅम्डा-परिभाष्य आहेत. तसेच \ref{lam:ldf:prf:lem:comp}
  आणि \ref{lam:ldf:prf:lem:prim} नुसार सर्वत्र परिभाषित लॅम्डा-परिभाष्य
  फलनांचा वर्ग संयोजन व आदिम पुनरावर्तन यांच्या
  अंतर्गत बंद आहे.
\end{proof}












\section{स्थिरबिंदू}\label{lam:ldf:fp:sec}

क्रमगुणित फलनाचे पुनरावर्तनाने असे पद $\fn{Fac}$
म्हणून परिभाषण करायचे आहे असे समजा:
\[\fn{Fac} \ident \lambd[n][\fn{IsZero}\, n\, \num 1 (\fn{Mult}\, n (\fn{Fac}(\fn{Pred} \, n)))]
\]
म्हणजे $n = 0$ असेल, तर $n$ चे क्रमगुणित $1$;
अन्यथा ते $n$ गुणिले $n-1$ चे क्रमगुणित.
पण $\fn{Fac}$ ची अशी व्याख्या करता येत नाही,
कारण उजव्या बाजूस $\fn{Fac}$ हेच पद येते.
अशा स्वसंदर्भाचा समावेश असलेल्या पुनरावर्ती
व्याख्या लॅम्डा कलनाचा भाग नाहीत. उदाहरणार्थ,
पुढील पद परिभाषित करताना
\[\fn{Mult} \ident \lambd[ab][b (\fn{Add}\, a) \num{0}]
\]
उजव्या बाजूस $\fn{Add}$ यांसारखी पूर्वी परिभाषित
पदेच वापरली जातात. $\fn{Add}$ च्या जागी त्याचे
परिभाषक पद ठेवून ते काढून टाकता येते. मग
$\fn{Mult}$ शुद्ध लॅम्डा पद म्हणून मिळेल;
$\fn{Add}$ मध्येही इतर परिभाषित पदे असतील
(उदा., $\fn{Add}'$ मध्ये आहेत), तर अशी जागापालट
सुरू ठेवून शेवटी शुद्ध लॅम्डा पद मिळवता येते.

वरच्या $\fn{Fac}$ सारख्या पुनरावर्ती व्याख्येच्या
बाबतीत मात्र तसे होत नाही. उजवीकडील $\fn{Fac}$
च्या जागी $\fn{Fac}$ चीच व्याख्या ठेवल्यास पुढील पद मिळते:
\begin{align*}
  B &\ident \lambd[m][\fn{IsZero}\,m\,\num{1}
       (\fn{Mult}\,m\,(\fn{Fac}(\fn{Pred}\,m)))]\\
  \fn{Fac} &\ident \lambd[n][\fn{IsZero}\,n\,\num{1}
       (\fn{Mult}\,n\,(B(\fn{Pred}\,n)))]
\end{align*}
येथे B हे फक्त आतले पद थोडक्यात लिहिण्यासाठीचे नाव आहे.
त्याच्या व्याख्येतील $\fn{Fac}$ अजूनही गेलेले नाही.
बाह्य अमूर्तीकरणाची व्याप्ती दोन्ही conditional शाखांवर आहे.
ही प्रक्रिया पुन्हा केल्यास व्याख्या वाढतच राहते
आणि कधीही शुद्ध लॅम्डा पद मिळत नाही. म्हणून
क्रमगुणित, किंवा सर्वसाधारणपणे पुनरावर्ती फलने,
अशा रीतीने परिभाषित करता येत नाहीत.

तरी या पुनरावर्ती व्याख्येतून एक गोष्ट समजते:
$f$ हे क्रमगुणित फलन दर्शवणारे पद असेल, तर
\[\fn{Fac}' \ident \lambd[g][\lambd[n][\fn{IsZero} \, n \, \num 1\, (\fn{Mult} \, n\, (g (\fn{Pred} n)))]]
\]
या पदाला $f$ लागू केल्यावर मिळणारे
$\fn{Fac}'\,f$ हेही क्रमगुणित फलन दर्शवते.
म्हणजे $\fn{Fac}'$ याकडे फलन स्वीकारून फलन
परत करणारे फलन म्हणून पाहिल्यास, $\fn{Fac'}\, f$
आणि~$f$ यांचे सर्व चर्च अंकांवरील परिणाम समान असतात.
यावरून ही दोन पदे स्वतः बीटा-सममूल्य आहेत असे मात्र
ठरत नाही; संख्यात्मक फलनाचे निरूपण आणि पदांची
बीटा-सममूल्यता या वेगळ्या अटी आहेत. $\fn{Fac}' \, f \equal[\beta] f$
हा गुणधर्म असलेल्या फलन~$f$ ला $\fn{Fac}'$
चा \emph{स्थिरबिंदू} म्हणतात. पुढील संयोजकाने
$\fn{Fac}'$ चा असा स्थिरबिंदू मिळवून त्याचे सर्व
अंकांवरील परिणाम क्रमगुणिताचे आहेत हे सिद्ध करू.

लॅम्डा कलनात दिलेल्या पदाचा स्थिरबिंदू काढणारी
पदे आहेत. त्यांच्या साहाय्याने $\fn{Fac}'$
सारख्या पदापासून क्रमगुणिताची व्याख्या मिळवता येते.

\begin{defn}
  \label{lam:ldf:fp:defn:Turing-Y}
  \emph{Y-संयोजक} हे पुढील पद आहे:
  \[  Y \ident (\lambd[ux][x(uux)])(\lambd[ux][x(uux)]).
  \]
\end{defn}

\begin{thm}
  $Y$ साठी कोणत्याही पद $g$ बाबत $Yg \red g(Yg)$; म्हणून
  $Yg$ हा नेहमी~$g$ चा स्थिरबिंदू असतो.
\end{thm}

\begin{proof}
  $(\lambd[ux][x(uux)])$ ला~$U$ असे संक्षेपाने
  लिहू. मग $Y \ident UU$. म्हणून
  \begin{align*}
    Y g &\ident (\lambd[ux][x(uux)])U\,g \\
    &\red (\lambd[x][x(UUx)])g\\
    & \red g(UUg) \ident g(Yg).
  \end{align*}
  $g(Yg)$ आणि $Yg$ या दोघांचेही $g(Yg)$ मध्ये
  न्यूनीकरण होते; म्हणून $g(Yg) \equal[\beta]
  Yg$ आणि $Yg$ हा~$g$ चा स्थिरबिंदू आहे.
\end{proof}

$Yg$ मध्ये रेडेक्स असल्याने न्यूनीकरण अनंतकाळ
पुढे चालू राहू शकते:
\begin{align*}
  Y g &\red g (Y g) \\
    &\red g (g (Y g)) \\
    &\red g(g (g (Y g)))\\
    &\ldots
\end{align*}
हा सांत पदांच्या अधिकाधिक उकलांचा अंतहीन अनुक्रम आहे;
तो स्वतः एक अनंत लॅम्डा पद नाही. म्हणून $Yg$ म्हणजे $g$
वारंवार लावलेल्या उकलांनी व्यक्त केलेला स्थिरबिंदू असे
सहज अर्थाने समजू शकतो. त्याला $g$ आणखी एकदा
लावल्याने, बोलायचे तर, काहीच जास्त होत नाही:
$Yg$ वर $g$ अनंत वेळा लावल्यावर पुन्हा $g$
लावले, तरी ते $Yg$ वर $g$ अनंत वेळा लावण्यासारखेच आहे.

वरील $Yg$ पासून सुरू होणारा $\beta$-न्यूनीकरणाचा
अनुक्रम अनंत आहे, हे लक्षात घ्या. त्यामुळे $Yg$
एखाद्या पदाला लागू केल्यावरचे $(Yg)N$ हे पदही
$(g(Yg))N$, $(g(g(Yg)))N$, \dots अशा अनंत
भिन्न पदांमध्ये न्यूनीत होऊ शकते. तरी न्यूनीकरणाचा
एखादा \emph{दुसरा} अनुक्रम सामान्य रूपात संपणे
शक्य आहे.

क्रमगुणिताचे उदाहरण घ्या. $\fn{Fac}$ ची व्याख्या
$Y\, \fn{Fac}'$ अशी करा (म्हणजे $\fn{Fac'}$
चा स्थिरबिंदू). मग:
\begin{align*}
  \fn{Fac}\,\num 3
  &\red Y\, \fn{Fac}' \, \num 3 \\
  &\red \fn{Fac}' (Y \,\fn{Fac}') \, \num 3 \\
  &\ident (\lambd[x][\lambd[n][\fn{IsZero} \, n\, \num 1\,
      (\fn{Mult}\, n\, (x (\fn{Pred}\, n)))]]) \, \fn{Fac} \, \num 3\\
  & \red \fn{IsZero} \, \num 3\, \num 1\,
      (\fn{Mult}\, \num 3\, (\fn{Fac} (\fn{Pred}\, \num 3))) \\
  &\red  \fn{Mult}\, \num 3 \, (\fn{Fac} \, \num 2).
  \intertext{त्याचप्रमाणे,}
  \fn{Fac}\,\num 2
  &\red  \fn{Mult}\, \num 2 \, (\fn{Fac} \, \num 1) \\
  \fn{Fac}\,\num 1
  &\red  \fn{Mult}\, \num 1 \, (\fn{Fac} \, \num 0)
  \intertext{पण}
  \fn{Fac}\,\num 0
  &\red \fn{Fac}' (Y \,\fn{Fac}') \, \num 0 \\
  &\ident (\lambd[x][\lambd[n][\fn{IsZero} \, n\, \num 1\,
      (\fn{Mult}\, n\, (x (\fn{Pred}\, n)))]]) \, \fn{Fac} \, \num 0\\
  & \red \fn{IsZero} \, \num 0\, \num 1\,
      (\fn{Mult}\, \num 0\, (\fn{Fac} (\fn{Pred}\, \num 0))).\\
  &\red \num 1.
  \intertext{म्हणून एकत्रितपणे}
  \fn{Fac}\,\num 3
  &\red  \fn{Mult}\, \num 3 \,
  (\fn{Mult} \, \num 2\, (\fn{Mult} \, \num 1\, \num 1)).
\end{align*}

$\fn{Fac'}$ साठी जे खरे आहे ते कोणत्याही
पुनरावर्ती व्याख्येसाठी खरे आहे. पुढील पुनरावर्ती
समीकरण दिले आहे असे समजा. पुनरावर्तनाचे नाव आणि
सर्व निविष्ट प्राचलांची नावे परस्पर भिन्न निवडायची आहेत:
\begin{align*}
g\,x_1\dots x_n & \equal[\beta] N
\intertext{येथे $N$ मध्ये~$g$ आणि $x_1$, \dots,~$x_n$
असू शकतात. मग असे पद~$G \ident (Y \lambd[g][\lambd[x_1\dots x_n][N]])$
नेहमी असते की}
G\,x_1 \dots x_n & \equal[\beta] \Subst{N}{G}{g}.
\intertext{कारण स्थिरबिंदू प्रमेयानुसार,}
G \ident (Y \lambd[g][\lambd[x_1\dots x_n][N]]) & \red \lambd[g][\lambd[x_1\dots x_n][N]](Y \lambd[g][\lambd[x_1\dots x_n][N]])\\
& \ident (\lambd[g][\lambd[x_1\dots x_n][N]])\,G
\intertext{आणि परिणामी,}
G\,x_1 \dots x_n
& \red (\lambd[g][\lambd[x_1\dots x_n][N]])\,G\,x_1\dots x_n\\
& \red (\lambd[x_1\dots x_n][\Subst{N}{G}{g}])\,x_1\dots x_n\\
  & \red \Subst{N}{G}{g}.
\end{align*}

\ref{lam:ldf:fp:defn:Turing-Y} मधील $Y$-संयोजक
ॲलन ट्यूरिंगने दिला होता. ॲलॉन्झो चर्चने
दुसरे रूप सुचवले होते; त्याला~$Y_C$ म्हणू:
\[Y_C \ident \lambd[g][(\lambd[x][g(xx)])(\lambd[x][g(xx)])].
\]
चर्चचा संयोजक ट्यूरिंगच्या संयोजकापेक्षा काहीसा
कमकुवत आहे. औपचारिक चर g वापरल्यास $Y_Cg \equal[\beta] g(Y_Cg)$, पण
$Y_Cg \bred g(Y_Cg)$ नाही. $V$ हे
$\lambd[x][g(xx)]$ पद असू द्या; येथे त्याचे बद्ध नाव
आदेशित पदाच्या मुक्त चरांपासून वेगळे निवडायचे आहे. मग
$Y_C \ident \lambd[g][VV]$. त्यामुळे
\begin{align*}
  VV & \ident (\lambd[x][g(xx)])V \red g(VV)
  \text{ आणि त्यामुळे}\\
  Y_C g & \ident (\lambd[g][VV])g \red VV \red g(VV), \text{ पण तसेच}\\
  g(Y_C g) & \ident g((\lambd[g][VV])g) \red g(VV).
\end{align*}
दुसऱ्या शब्दांत, $Y_Cg$ आणि $g(Y_Cg)$ यांचे
न्यूनीकरण होऊन एकच पद $g(VV)$ मिळते;
म्हणून $Y_Cg \equal[\beta] g(Y_Cg)$. अनेक
उपयोजनांसाठी एवढे पुरेसे असते.











\section{लघुतमीकरण}\label{lam:ldf:min:sec}

मूलभूत फलने $\Zero$, $\Succ$, $\Proj{n}{i}$
यांपासून संयोजन, आदिम पुनरावर्तन आणि नियमित
लघुतमीकरण वापरून मिळणारी फलने म्हणजे सामान्य
पुनरावर्ती फलने. सर्व सामान्य पुनरावर्ती फलने
लॅम्डा-परिभाष्य आहेत हे दाखवण्यासाठी,
लॅम्डा-परिभाष्य फलनापासून नियमित लघुतमीकरणाने
परिभाषित केलेले कोणतेही फलनही लॅम्डा-परिभाष्य
आहे हे दाखवावे लागेल.

\begin{lem}
  \label{lam:ldf:min:lem:min} $f(x_1, \dots, x_k, y)$ हे सर्वत्र परिभाषित आणि नियमित
  आणि लॅम्डा-परिभाष्य असेल, तर
  \[  h(x_1, \dots, x_k) = \umin{y}{f(x_1,\dots,x_k, y) = 0}
  \]
  याने परिभाषित केलेले $h$~देखील
  लॅम्डा-परिभाष्य आहे.
\end{lem}

\begin{proof}
  लॅम्डा पद~$F$ हे नियमित फलन $f(\vec x, y)$ ला
  $\lambda$-परिभाषित करते असे समजा. $h$ ला
  लॅम्डा-परिभाषित करण्यासाठी शोधफलन आणि
  स्थिरबिंदू संयोजक वापरू:
  \begin{align*}
    \fn{Search} & \ident \lambd[g][\lambd[f\,\vec{x}\,y][
        \fn{IsZero} (f\, \vec{x}\, y)\, y\, (g\, f\, \vec{x} (\fn{Succ}\, y))]]\\
    H & \ident \lambd[\vec x][(Y \, \fn{Search}) F\, \vec{x}\, \num{0}],
  \end{align*}
  येथे $Y$ हा \ref{lam:ldf:fp:defn:Turing-Y} मधील
  ट्यूरिंग स्थिरबिंदू संयोजक आहे. Recursive call मध्ये
  predicate प्रतिनिधी पुन्हा पहिला फलसाधक म्हणून द्यायचा
  आहे; तो वगळल्यास पुढील calls चे सर्व फलसाधक सरकतात.
  सहज अर्थाने, $\fn{Search}$ हे स्वतःला पुन्हा
  बोलावणारे फलन आहे: $y$ पासून सुरू करून
  $f\, \vec x\, y$ शून्य आहे का हे तपासते;
  शून्य असेल तर~$y$ परत करते, अन्यथा
  $\fn{Succ}\, y$ घेऊन स्वतःला पुन्हा बोलावते.
  म्हणून $(Y \, \fn{Search}) F
  \num{n_1}\dots\num{n_k}\,\num{0}$ हे
  $f(n_1, \dots, n_k, m) = 0$ पूर्ण करणारा
  लघुतम~$m$ परत करते.

  ट्यूरिंग संयोजकाच्या थेट उकलीने आणि परिभाषित
  प्रत्येक predicate चाचणीच्या न्यूनीकरणाने पुढील गोष्ट मिळते:
  \begin{align*}
    (Y \, \fn{Search}) F \num{n_1}
    \dots\num{n_k}\, \num{m} & \red \num{m}
    \intertext{जर $f(n_1, \dots,
      n_k, m) = 0$ असेल; अन्यथा}
    & \red (Y \, \fn{Search}) F\, \num{n_1} \dots\num{n_k}\, \num{m+1}
    \intertext{मिळते. $f$ नियमित असल्याने काही $y$ साठी
      $f(n_1, \dots, n_k, y) = 0$ असते; म्हणून}
    (Y \, \fn{Search}) F \num{n_1} \dots\num{n_k}\,\num{0}
    & \red \num{h(n_1, \dots, n_k)}.
    \end{align*}
\end{proof}

\begin{prop}
  प्रत्येक सामान्य पुनरावर्ती फलन लॅम्डा-परिभाष्य आहे.
\end{prop}

\begin{proof}
 \ref{lam:ldf:prf:lem:basic} नुसार सर्व मूलभूत फलने
 लॅम्डा-परिभाष्य आहेत. तसेच \ref{lam:ldf:prf:lem:comp},
 \ref{lam:ldf:prf:lem:prim} आणि \ref{lam:ldf:min:lem:min} नुसार
 सर्वत्र परिभाषित लॅम्डा-परिभाष्य फलनांचा वर्ग संयोजन,
 आदिम पुनरावर्तन आणि नियमित लघुतमीकरण यांच्या
 अंतर्गत बंद असतो.
\end{proof}











\section{आंशिक पुनरावर्ती फलने लॅम्डा-परिभाष्य असतात}\label{lam:ldf:par:sec}

मूलभूत फलनांपासून संयोजन, आदिम पुनरावर्तन आणि
अपरिबद्ध लघुतमीकरण वापरून मिळणारी फलने म्हणजे
आंशिक पुनरावर्ती फलने. सामान्य पुनरावर्ती फलनांपेक्षा
त्यांचा फरक असा की अपरिबद्ध शोधात वापरलेली फलने
नियमित असणे आवश्यक नाही. नियमिततेची अट नसल्याने
लघुतमीकरणाने परिभाषित केलेली फलने काही वेळा
अपरिभाषित राहू शकतात.

प्रथमदर्शनी, (सर्वत्र परिभाषित) सामान्य पुनरावर्ती
फलने लॅम्डा-परिभाष्य आहेत हे दाखवण्याच्या
पद्धतीनेच सर्व आंशिक पुनरावर्ती फलनेही
लॅम्डा-परिभाष्य आहेत हे सिद्ध करता येईल असे
वाटू शकते. उदाहरणार्थ, $f$ आणि $g$ ही अनुक्रमे
$F$ आणि~$G$ या पदांनी लॅम्डा-परिभाषित
असतील, तर $f$ चे $g$ शी संयोजन
$\lambd[x][F (G x)]$ ने लॅम्डा-परिभाषित होते,
जर दोन्ही फलने सर्वत्र परिभाषित असतील.
पण फलने आंशिक असतील, तेव्हा अडचण येते.
$g(x)$ अपरिभाषित असल्यास $G x$ ला सामान्य रूप
नसते. बहुतेक वेळा मग $F (G x)$ लाही सामान्य रूप
नसते, आणि आपल्याला तेच हवे आहे. पण
$F$ हे $\lambd[x][\num{0}]$ असेल, तर
$F (G x)$ ला $\num{0}$ हे सामान्य रूप मिळते.

ही अडचण दूर करता येते. सर्व आंशिक पुनरावर्ती
फलने अशी लॅम्डा-परिभाषित करण्याचे मार्ग आहेत
की अपरिभाषित मूल्ये सामान्य रूप नसलेल्या पदांनी
दर्शवली जातात. मात्र हे मार्ग सामान्य पुनरावर्ती
फलनांसाठी वापरलेल्या पद्धतीपेक्षा काहीसे अधिक
गुंतागुंतीचे आणि कमी सहज आहेत. येथे पुरावा न देता
प्रमेय नोंदवतो. या आवृत्तीच्या आधीच्या परिचयात
\ref{lam:int:lrp:thm:computable-lambda} ची दुरुस्त
सिद्धता आणि लघुतमीकरणातील forced-search extraction
यांनी ही construction आधीच पूर्ण केली आहे:

\begin{thm}
  सर्व आंशिक पुनरावर्ती फलने लॅम्डा-परिभाष्य आहेत.
\end{thm}











\section{लॅम्डा-परिभाष्य फलने पुनरावर्ती असतात}\label{lam:dfl:ldr:sec}

सर्व आंशिक पुनरावर्ती फलने लॅम्डा-परिभाष्य
आहेत एवढेच नव्हे; याचे व्यस्त विधानही खरे आहे.
म्हणजे सर्व लॅम्डा-परिभाष्य फलने आंशिक
पुनरावर्ती असतात.

\begin{thm}
  \label{lam:dfl:ldr:thm:lambda-computable} आंशिक फलन $f$
  हे लॅम्डा-परिभाष्य असेल, तर ते आंशिक
  पुनरावर्ती असते.
\end{thm}

\begin{proof}
  पुराव्याची फक्त रूपरेषा देऊ. आधी $\lambd$-पदांचे
  अंकगणितीकरण करू: अनुक्रमांच्या नेहमीच्या अभाज्य
  संख्यांच्या घातांवरील संकेतनाचा वापर करून
  $\lambd$-पदांना पद्धतशीरपणे गोडेल संख्या देऊ.
  मग आंशिक पुनरावर्ती फलन $\fn{normalize}(t)$
  परिभाषित करू. ते लॅम्डा पदाची गोडेल संख्या~$t$
  फलसाधक म्हणून घेते आणि पदाला सामान्य रूप
  असेल, तर त्या रूपाची गोडेल संख्या परत करते;
  अन्यथा ते अपरिभाषित राहते. हे फलन प्रभावी आहे:
  प्रत्येक पायरीत डावीकडील सर्वप्रथम बीटा रेडेक्स
  निवडून, capture टाळण्यासाठी ठरावीक क्रमाने ताजी
  बद्ध नावे घेऊन, त्याचे संकुचन करायचे. सामान्य रूप
  असल्यास ही पद्धत तिथे पोहोचते; अन्यथा शोध संपत नाही.
  अवैध input codes वरही फलन अपरिभाषित ठेवू. पुढे
  $\fn{toChurch}$ आणि $\fn{fromChurch}$ ही दोन
  आंशिक पुनरावर्ती फलने परिभाषित करू. ती नैसर्गिक
  संख्या आणि त्या संख्यांच्या अनुरूप चर्च अंकांच्या
  गोडेल संख्या यांमध्ये पुढे व मागे रूपांतर करतात.
  पहिले रूपांतर सर्व नैसर्गिक संख्यांवर परिभाषित आहे.
  दुसरे रूपांतर फक्त चर्च अंकाच्या सामान्य रूपाचा वैध
  code दिल्यास परिभाषित आहे; alpha-सममूल्य बद्ध
  नावांची निवड बदलली, तरी अंक तोच ओळखायचा.

  या पुनरावर्ती फलनांचा वापर करून फलन~$f$
  आंशिक पुनरावर्ती म्हणून परिभाषित करता येते.
  $f$ ला लॅम्डा-परिभाषित करणारे $\lambd$-पद~$F$
  आहे. $f(n_1, \dots, n_k)$ संगणित करण्यासाठी
  आधी $\fn{toChurch}(n_i)$ वापरून अनुरूप चर्च
  अंकांच्या गोडेल संख्या मिळवा. $\Gn{F}$ आणि या
  codes पासून संकेतनातील प्रभावी उपयोजन रचनाकार
  वापरून $F \num{n_1}\dots\num{n_k}$ या पदाची
  गोडेल संख्या मिळवा; येथे साधी numeric बेरीज किंवा
  codes चे अंध जोडणे अभिप्रेत नाही. आता या गोडेल संख्येवर
  $\fn{normalize}$ वापरा. $f(n_1, \dots, n_k)$
  परिभाषित असेल, तर
  $F \num{n_1}\dots\num{n_k}$ ला सामान्य रूप
  असते (ते चर्च अंकच असले पाहिजे); अन्यथा सामान्य
  रूप नसते (आणि म्हणून
  \[\fn{normalize}(\Gn{F\num{n_1}\dots\num{n_k}})\]
  हे अपरिभाषित असते). शेवटी सामान्यीकृत पदाच्या
  गोडेल संख्येवर $\fn{fromChurch}$ वापरा.
\end{proof}



\part{बहुमूल्य तर्कशास्त्र}\label{mvl:part}
\begin{editorial}
  या भागात विधानीय बहुमूल्य तर्कशास्त्रांवरील मसुदा साहित्य आहे.
\end{editorial}
